\section{Discussion}
\label{sec:discussion}

\textbf{What remains for devices away from home.} DriftGate gains most where the device exit is strong and gives up a little where the edge exit is strong. A rule that knew the kind of every request, and answered Main requests with the device exit and all other requests with the edge exit, would reach 72.9\% in S1, compared with 68.3\% for DriftGate. Most of this gap lies in the OOP and OOR requests of devices away from home. Closing it requires a signal that recognizes, for a single request, that its class is new to the device. The confidence of the device exit does not provide such a signal (Section~\ref{sec:choose}), and label-free features of both exits were not enough even when a classifier was trained on them with labels (Section~\ref{sec:overall}). Distances in the feature space of the device block, measured against the device's own training data, are a natural next candidate.

\textbf{Relation to training.} DriftGate acts after training and only on the answers of the requesting device. It therefore avoids the trade-off of Section~\ref{sec:training_side}, where a device that personalizes more helps itself at the expense of visitors to its cell. Because DriftGate needs only the two output vectors, the class list of the device, and the class counts of the cell, it can run on top of other ways of training the two exits, such as different personalization ratios or personalized federated learning objectives.

\textbf{Information at the edge.} The prior correction uses the device's class list $M_k$ and the class counts of the cell. In split learning, the edge receives training labels and already holds both. In variants that keep labels on the device, the device can send its class list once, which reveals no more than which classes it trains on.

\textbf{Scope.} We studied changes in which classes a device sees, which is how mobility changes the traffic of a recognition application. Changes in how the same classes look, such as lighting or sensor differences between places, affect both exits and are outside the scope of the correction. Our scenarios train the models from random initialization over one simulated day, so the first hours of each run include early training. In longer deployments, both exits are already trained when a day starts, and we expect the morning hours to behave like the evening hours of our runs, when every device is again at home and both exits are trained.
